\subsection{双自变量函数}

二元函数是指其定义域为有序对集合（或等价地，为某个乘积的子集）的函数。设 \(f\) 是这样一个函数；如果 \((x, y)\) 是定义域中的一个元素 \(f\) ，其值 \(f((x, y))\) 的 \(f\) 在元素 \((x, y)\) 通常用 \(f(x, y)\) . 设 \(f\) 是两个变量的函数， \(D\) 其定义域，以及 \(C\) 其陪域。对于每个 \(y\) 让 \(A_y\) 是所有……的集合 \(x\) 使得 \((x, y) \in D\) （即 \(\overline{D}^1\) 在 \(y\) （第1号）的截面）。映射

\[
x \rightarrow f (x, y) \quad (x \in \mathrm{A} _ {y}, f (x, y) \in \mathrm{C})
\]

称为由 \(f\) 关于第二个参数的取值 \(y\) 所定义的偏映射，并记为 \(f(\bullet, y)\) ，或 \(f(., y)\) 的复合（有时也称为 \(f_y\) ）；我们有 \(f(\bullet, y)(x) = f(x, y)\) 对于所有 \((x, y) \in D\) . 类似地，对于每一个 \(x\) 让 \(B_x\) 是所有……的集合 \(y\) 使得 \((x, y) \in D\) . 映射

\[
y \rightarrow f (x, y) \quad (y \in \mathbf {B} _ {x}, f (x, y) \in \mathbf {C})
\]

称为由 \(f\) 关于第二个参数的取值 \(x\) 关于第一个参数的取值，并记为 \(f(x, \bullet)\) ，或 \(f(x, \cdot)\) 的复合（有时也称为 \(f_x\) ）；我们有 \(f(x, \bullet)(y) = f(x, y)\) 对于所有 \((x, y) \in \mathbf{D}\) 。如果对每个 \(y\) （分别） \(x\) )，偏映射 \(f(\bullet, y)\) （分别） \(f(x, \bullet)\) )是常值映射，则称 \(f\) 不依赖于其第一个（分别地，第二个）参数；这意味着因此 \(f(x, y) = f(x', y)\) 每当 \((x, y)\) 并且 \((x', y)\) 属于 \(\mathbf{D}\) （分别） \(f(x, y) = f(x, y')\) 每当 \((x, y)\) 并且 \((x, y')\) 属于 \(\mathbf{D}\) ). 对于每一个 \(y\) 属于 \(\mathbf{D}\) 让 \(g(y)\) 表示其公共值 \(f(x, y)\) 为了 \(x \in A_y\) ；该映射 \(y \to g(y)\) 那么就是一个映射， \(\mathrm{pr}_2\mathbf{D}\) 进入。 \(\mathbf{C}\) 使得

\[
g (y) = f (x, y) \quad \text {   对   所有   } (x, y) \in D.
\]

反之，设 \(g\) 是集合 \(B\) 到集合 \(C\) 的一个映射，并设 \(A\) 为任意集合。则映射 \((x, y) \to g(y)\) （从 \(A \times B\) 到 \(C\) ）不依赖于第一个参数。

让 \(u\) 是A到C的一个映射，而 \(v\) B到D的一个映射。映射 \(z \to (u(\mathrm{pr}_1z), v(\mathrm{pr}_2z))\) 的 \(A \times B\) 进入。 \(C \times D\) 被称为 \(u\) 并且 \(v\) 到乘积集的（典范）扩张，或简称为 \(u\) 并且 \(v\) （若无混淆风险）；其值域为 \(u\langle A\rangle \times v\langle B\rangle\) 。它记作 \(u \times v\) 。如果 \(u\) 并且 \(v\) 是单射（相应地，满射），则 \(u \times v\) 是单射（相应地，满射）。如果 \(u\) 并且 \(v\) 若它们是双射的，则 \(u \times v\) 是双射的，且其逆映射为 \(\overline{u}^1 \times \overline{v}^1\) 。如果 \(u'\) 是将C映射到集合E的映射，并且如果 \(v'\) 是D到集合F的一个映射，我们有

\[
\left(u ^ {\prime} \times v ^ {\prime}\right) \circ (u \times v) = \left(u ^ {\prime} \circ u\right) \times \left(v ^ {\prime} \circ v\right).
\]

如果U、V分别是u、v的图像，那么W的图像 \(u \times v\) 是序对的集合 \(((x, y), (z, t))\) 的 \((\mathbf{A} \times \mathbf{B}) \times (\mathbf{C} \times \mathbf{D})\) 使得 \((x, z) \in \mathbf{U}\) 并且 \((y, t) \in \mathbf{V}\) ；该映射 \(((x, y), (z, t)) \to ((x, z), (y, t))\) 将W与乘积建立一一对应关系 \(U \times V\) （的一个子集 \((\mathbf{A} \times \mathbf{C}) \times (\mathbf{B} \times \mathbf{D})\) （参见§5，第5节）。

